\documentclass[12pt,a4paper]{article}

\usepackage[margin=2.5cm]{geometry}
\usepackage{booktabs}
\usepackage{array}
\usepackage{amsmath}
\usepackage{amsfonts}
\usepackage{listings}
\usepackage{xcolor}
\usepackage{hyperref}
\usepackage{caption}
\usepackage{float}
\usepackage{graphicx}
\usepackage{parskip}
\usepackage{longtable}

\hypersetup{
    colorlinks=true,
    linkcolor=black,
    urlcolor=blue,
    citecolor=black
}

\lstset{
    basicstyle=\ttfamily\small,
    keywordstyle=\color{blue},
    commentstyle=\color{gray},
    stringstyle=\color{orange},
    frame=single,
    breaklines=true,
    numbers=left,
    numberstyle=\tiny\color{gray},
    captionpos=b
}

\title{
    \textbf{COS314 : Artificial Intelligence} \\[0.5em]
    \large Assignment Three: Genetic Programming for Breast Cancer Classification \\[0.3em]
    \normalsize Department of Computer Science \\
    Engineering, Built Environment and Information Technology \\
    University of Pretoria
}
\author{Finnley Wyllie \and Alessandro Paravano}
\date{2026}

\begin{document}

\maketitle
\tableofcontents
\newpage

% -----------------------------------------------------------------------
\section{Introduction}
% -----------------------------------------------------------------------

This report describes the design, implementation, and evaluation of two Genetic Programming (GP) systems developed to classify the Breast Cancer Wisconsin dataset into \emph{no-recurrence-events} (class 0) and \emph{recurrence-events} (class 1). The two approaches are:

\begin{enumerate}
    \item \textbf{Arithmetic GP (Symbolic Classifier)} : evolves arithmetic expressions over the feature variables whose sign determines the predicted class.
    \item \textbf{Decision Tree GP (Logical Classifier)} : evolves binary decision trees whose branching structure classifies instances using threshold comparisons on individual features.
\end{enumerate}

Both systems are implemented in Java, seeded for full reproducibility, and evaluated over 30 independent runs. The best-performing run from each approach is reported and compared using a statistical significance test.

% -----------------------------------------------------------------------
\section{Dataset}
% -----------------------------------------------------------------------

\subsection{Description}

The Breast Cancer Wisconsin (Diagnostic) dataset contains 286 instances. A training/test split is provided via \texttt{Breast\_train.csv} and \texttt{Breast\_test.csv}. Each instance is described by nine features, encoded as integers according to the mappings in Table~\ref{tab:features}.

\begin{table}[H]
\centering
\caption{Feature encoding summary}
\label{tab:features}
\begin{tabular}{lll}
\toprule
\textbf{Feature} & \textbf{Original Values} & \textbf{Encoded Values} \\
\midrule
Age              & 20--29, 30--39, \ldots, 70--79 & 0--5 \\
Menopause        & premeno, ge40, lt40            & 0, 1, 2 \\
Tumor Size       & Binned ranges                  & Integer mapping \\
Inv-Nodes        & Binned ranges                  & Integer mapping \\
Node-Caps        & no, yes, ?                     & 0, 1, 2 \\
Deg-Malig        & 1, 2, 3                        & 1, 2, 3 \\
Breast           & left, right                    & 0, 1 \\
Breast-Quad      & left-low, right-up, \ldots     & 0--5 \\
Irradiat         & no, yes, ?                     & 0, 1, 2 \\
\midrule
Class            & no-recurrence, recurrence      & 0, 1 \\
\bottomrule
\end{tabular}
\end{table}

\subsection{Pre-processing}

No further normalisation or imputation was applied beyond the integer encoding provided in the dataset files. All features are read directly as floating-point values. The class label is read from the first column of each CSV row, and the remaining nine columns form the feature vector.

% -----------------------------------------------------------------------
\section{Methodology}
% -----------------------------------------------------------------------

\subsection{Shared GP Framework}

Both classifiers share a common abstract base class (\texttt{GeneticProgram}) that implements:

\begin{itemize}
    \item \textbf{Ramped half-and-half initialisation} : the population is divided into equal segments, one per depth level from 1 to \texttt{maxDepth}. Within each segment, alternating individuals are generated using the \emph{full} and \emph{grow} methods, ensuring diversity in tree shape and size.
    \item \textbf{Tournament selection} : $k$ individuals are drawn at random and the fittest is selected as a parent.
    \item \textbf{Generational replacement} : the entire population is replaced by offspring each generation, with elitism achieved by tracking the all-time best individual separately.
    \item \textbf{Early termination} : evolution halts immediately if a perfect fitness of 1.0 is reached.
\end{itemize}

A seeded \texttt{java.util.Random} instance is used for all stochastic decisions, ensuring full reproducibility.

\subsection{Arithmetic GP (Symbolic Classifier)}

\subsubsection{Representation}

Each individual is a binary expression tree. Internal nodes are arithmetic operators; leaf nodes are either a named feature variable or a random constant drawn uniformly from $[-1, 1]$.

\textbf{Function set:} $\{+,\;-,\;\times,\;\div\}$

Protected division is used: if the denominator is zero, the result is 1.0.

\textbf{Terminal set:} $\{\texttt{age},\;\texttt{menopause},\;\texttt{tumor\_size},\;\texttt{inv\_nodes},\;\texttt{node\_caps},\;\texttt{deg\_malig},\;\texttt{breast},\;\texttt{breast\_quad},\;\texttt{irradiat}\} \cup \mathbb{R}_{[-1,1]}$

\subsubsection{Classification Rule}

\[
\hat{y} =
\begin{cases}
1 & \text{if } f(\mathbf{x}) \geq 0 \\
0 & \text{otherwise}
\end{cases}
\]

where $f(\mathbf{x})$ is the arithmetic expression evaluated with the feature values substituted for the variable nodes.

\subsubsection{Crossover}

Subtree crossover is applied. A random \texttt{FunctionNode} is selected in each parent tree, and one of its subtrees (left or right, chosen uniformly) is swapped between the two children. The resulting trees are pruned back to \texttt{maxDepth} by replacing any over-deep node's children with random terminals.

\subsubsection{Mutation}

Point mutation is applied with equal probability between two variants:
\begin{itemize}
    \item \textbf{Grow mutation} : a random function node is selected and one of its children is replaced by a newly generated subtree of depth up to \texttt{mutationDepth}.
    \item \textbf{Shrink mutation} : a random function node is selected and one of its children is replaced by a random terminal node.
\end{itemize}

After grow mutation the tree is pruned to enforce \texttt{maxDepth}.

% -----------------------------------------------------------------------
\subsection{Decision Tree GP (Logical Classifier)}

\subsubsection{Representation}

Each individual is a binary decision tree. Internal nodes are \texttt{DTDecisionNode} objects that encode a threshold comparison of the form $\texttt{feature} < \texttt{threshold}$; leaf nodes are \texttt{DTLeafNode} objects that carry a class label of 0 or 1.

\textbf{Decision node:} \texttt{IF(feature $<$ threshold, left\_subtree, right\_subtree)}

Thresholds are drawn uniformly from $[0, 11]$, covering the full range of encoded feature values.

\subsubsection{Classification Rule}

Tree traversal: at each internal node, if the named feature value is less than the threshold the left branch is taken; otherwise the right branch is taken. The class label at the reached leaf is the prediction.

\[
\hat{y} =
\begin{cases}
1 & \text{if } \mathrm{evaluate}(\text{tree}) \geq 0.5 \\
0 & \text{otherwise}
\end{cases}
\]

\subsubsection{Crossover}

Subtree crossover between two decision trees: a random \texttt{DTDecisionNode} is selected in each parent, a subtree is extracted from each, and the subtrees are swapped. Children are pruned to \texttt{maxDepth}.

\subsubsection{Mutation}

Point mutation selects any node uniformly at random:
\begin{itemize}
    \item \textbf{Decision node} : with equal probability, either the threshold is replaced by a new random value in $[0, 11]$, or the feature name is replaced by a random feature from the feature set.
    \item \textbf{Leaf node} : the class label is flipped ($0 \leftrightarrow 1$).
\end{itemize}

% -----------------------------------------------------------------------
\section{Experimental Setup}
% -----------------------------------------------------------------------

\subsection{Parameters}

Table~\ref{tab:params} lists all GP parameters used for both models.

\begin{table}[H]
\centering
\caption{Genetic Programming parameters}
\label{tab:params}
\begin{tabular}{ll}
\toprule
\textbf{Parameter}          & \textbf{Value} \\
\midrule
Population size             & 200 \\
Initial tree generation     & Ramped half-and-half \\
Initial tree depth          & 1 -- 5 \\
Max offspring depth         & 5 \\
Selection method            & Tournament selection \\
Tournament size             & 5 \\
Function set (Arithmetic)   & $\{+,\,-,\,\times,\,\div\}$ \\
Function set (Decision Tree) & Threshold comparisons \\
Crossover rate              & 80\% \\
Mutation rate               & 20\% \\
Mutation type               & Point (grow / shrink / flip) \\
Mutation offspring depth    & 3 \\
Fitness function            & Classification accuracy \\
Maximum generations         & 100 \\
\bottomrule
\end{tabular}
\end{table}

\subsection{Reproducibility}

The program prompts for a seed value before execution. The seed is passed directly to \texttt{java.util.Random} and governs all stochastic operations. The seed used for the best run of each model is reported in Section~\ref{sec:results} so that results can be independently replicated.

\begin{lstlisting}[language=Java, caption={Seeded random initialisation (GeneticProgram.java)}]
public GeneticProgram(..., long seed) {
    ...
    this.random = new Random(seed);
}
\end{lstlisting}

\subsection{Execution}

The program is compiled and run using:

\begin{lstlisting}[language=bash]
make        # compiles all .java files
make run    # compiles and launches Main
\end{lstlisting}

At startup the user is prompted for: model choice, mode (demonstration / classification), seed value, training file path, and test file path.

\subsection{Evaluation Protocol}

Both models were evaluated over \textbf{30 independent runs}, each with a unique seed. The best-performing run (highest test accuracy) is reported as the final result and will be demonstrated on demo day. Performance is measured by:

\begin{itemize}
    \item \textbf{Training accuracy} : proportion of correctly classified training instances.
    \item \textbf{Test accuracy} : proportion of correctly classified unseen test instances.
    \item \textbf{F-measure} : harmonic mean of precision and recall on the test set:
    \[
        F_1 = \frac{2 \cdot \mathrm{Precision} \cdot \mathrm{Recall}}{\mathrm{Precision} + \mathrm{Recall}}
    \]
    \item \textbf{Runtime} : wall-clock time from start of training to end of training, in milliseconds.
\end{itemize}

% -----------------------------------------------------------------------
\section{Results}
\label{sec:results}
% -----------------------------------------------------------------------

\subsection{Best-Run Results}

Table~\ref{tab:results} presents the metrics for the best run of each model.

\begin{table}[H]
\centering
\caption{Comparison of classification performance (best run from 30 independent runs)}
\label{tab:results}
\begin{tabular}{lllll}
\toprule
\textbf{Algorithm}   & \textbf{Training (\%)} & \textbf{Test (\%)} & \textbf{F-measure} & \textbf{Runtime (ms)} \\
\midrule
Arithmetic GP        & 83.0601                & 62.7907            & 0.5556             & 523 \\
Decision Tree GP     & 84.1530                & 58.1395            & 0.4000             & 463 \\
\bottomrule
\end{tabular}
\end{table}

\noindent\textbf{Best seed (Arithmetic GP):} \texttt{25}

\noindent\textbf{Best seed (Decision Tree GP):} \texttt{1}

\subsection{30-Run Summary}

Table~\ref{tab:runs} summarises the distribution of test accuracy across all 30 runs.

\begin{table}[H]
\centering
\caption{Summary statistics over 30 independent runs}
\label{tab:runs}
\begin{tabular}{lllll}
\toprule
\textbf{Algorithm}  & \textbf{Mean (\%)} & \textbf{Std Dev} & \textbf{Min (\%)} & \textbf{Max (\%)} \\
\midrule
Arithmetic GP       & 47.48 & 5.27 & 39.53 & 62.79 \\
Decision Tree GP    & 50.50 & 3.65 & 43.02 & 58.14 \\
\bottomrule
\end{tabular}
\end{table}

\subsection{30-Run Results Table}

Table~\ref{tab:runresults} lists the individual evaluation results for all 30 independent runs of the Arithmetic GP. These are the seeded runs used to determine the best-performing model for demonstration.

\small
\begin{longtable}{rrrrr}
\caption{Per-run evaluation results for 30 independent Arithmetic GP runs}\label{tab:runresults}\\
\toprule
\textbf{Seed} & \textbf{Train (\%)} & \textbf{Test (\%)} & \textbf{F-measure} & \textbf{Runtime (ms)} \\
\midrule
\endfirsthead

\toprule
\textbf{Seed} & \textbf{Train (\%)} & \textbf{Test (\%)} & \textbf{F-measure} & \textbf{Runtime (ms)} \\
\midrule
\endhead

1  & 82.5137 & 47.6744 & 0.4000 & 587 \\
2  & 83.0601 & 48.8372 & 0.4634 & 418 \\
3  & 83.0601 & 47.6744 & 0.2623 & 584 \\
4  & 84.6995 & 45.3488 & 0.3896 & 641 \\
5  & 83.6066 & 43.0233 & 0.2462 & 483 \\
6  & 84.1530 & 39.5349 & 0.2571 & 558 \\
7  & 83.0601 & 47.6744 & 0.4000 & 615 \\
8  & 83.6066 & 48.8372 & 0.3889 & 565 \\
9  & 83.6066 & 45.3488 & 0.4337 & 550 \\
10 & 84.6995 & 39.5349 & 0.3333 & 653 \\
11 & 81.4208 & 50.0000 & 0.2951 & 383 \\
12 & 84.1530 & 50.0000 & 0.4819 & 549 \\
13 & 84.6995 & 39.5349 & 0.2973 & 793 \\
14 & 83.0601 & 52.3256 & 0.4675 & 602 \\
15 & 83.6066 & 50.0000 & 0.4941 & 573 \\
16 & 84.1530 & 40.6977 & 0.2609 & 588 \\
17 & 84.6995 & 52.3256 & 0.3492 & 596 \\
18 & 84.6995 & 43.0233 & 0.2899 & 605 \\
19 & 84.1530 & 43.0233 & 0.3099 & 812 \\
20 & 83.6066 & 45.3488 & 0.3380 & 661 \\
21 & 83.0601 & 56.9767 & 0.5316 & 470 \\
22 & 84.1530 & 44.1860 & 0.2941 & 609 \\
23 & 81.9672 & 52.3256 & 0.4225 & 594 \\
24 & 84.1530 & 47.6744 & 0.3662 & 528 \\
25 & 83.0601 & 62.7907 & 0.5556 & 523 \\
26 & 83.0601 & 46.5116 & 0.3429 & 623 \\
27 & 83.0601 & 43.0233 & 0.3467 & 532 \\
28 & 84.1530 & 52.3256 & 0.3279 & 658 \\
29 & 82.5137 & 52.3256 & 0.3051 & 830 \\
30 & 83.6066 & 46.5116 & 0.4390 & 599 \\
\bottomrule
\end{longtable}
\normalsize

% -----------------------------------------------------------------------
\section{Statistical Significance}
% -----------------------------------------------------------------------

A \textbf{Wilcoxon signed-rank test} was applied to the test accuracy values from all 30 independent runs of each model to determine whether the observed difference in performance is statistically significant (significance level $\alpha = 0.05$).

The Wilcoxon test is preferred over a paired $t$-test here because classification accuracy is bounded $[0, 1]$ and may not follow a normal distribution across runs.

\vspace{1em}
\noindent\textbf{Null hypothesis} $H_0$: There is no statistically significant difference in classification performance between the Arithmetic GP and the Decision Tree GP.

\vspace{0.5em}
\noindent\textbf{Result:} $W = 80.0$, $z = -2.619$, $p = 0.0088$

\vspace{0.5em}
\noindent\textbf{Conclusion:} We reject $H_0$ at $\alpha = 0.05$. The Decision Tree GP achieves a statistically significantly higher mean test accuracy (50.50\%) than the Arithmetic GP (47.48\%) across 30 independent runs, despite the Arithmetic GP producing the single best-run score of 62.79\%.

% -----------------------------------------------------------------------
\section{Discussion}
% -----------------------------------------------------------------------

\subsection{Arithmetic GP}

The arithmetic classifier evolves symbolic expressions directly in the feature space. The sign of the expression value is the decision boundary, making it interpretable as a hyperplane in the transformed feature space. Protected division prevents division-by-zero failures during evaluation.

Mutation uses two complementary operators: grow mutation increases complexity by inserting subtrees, while shrink mutation reduces complexity by collapsing subtrees to terminals. Alternating between them helps balance exploration and exploitation.

\subsection{Decision Tree GP}

The decision tree GP evolves explicit if-then-else structures over feature thresholds, which are inherently interpretable. Point mutation directly modifies either the split feature or the split threshold at a single node, or flips a leaf label, providing fine-grained local search.

The maximum threshold of 11.0 was chosen to cover the widest observed encoded feature value across all nine attributes, ensuring all possible splits are reachable.

\subsection{Comparison}

Across 30 independent runs the Decision Tree GP was more \textbf{consistent}, with a lower standard deviation of 3.65\% versus 5.27\% for the Arithmetic GP. The Arithmetic GP exhibited higher variance, occasionally reaching a peak test accuracy of 62.7907\% (seed 25) but also falling as low as 39.5349\%. Both models show signs of mild overfitting, with training accuracies around 83--84\% against test accuracies near 47--50\% on average, reflecting the limited and imbalanced nature of the dataset.

The Decision Tree GP was also faster at 463 ms for its best run versus 523 ms, owing to its simpler tree evaluation (threshold comparisons versus floating-point arithmetic). The Wilcoxon signed-rank test ($p = 0.0088$) confirms the Decision Tree GP's advantage over 30 runs is statistically significant, even though a single lucky Arithmetic GP run outperformed the best Decision Tree result.

% -----------------------------------------------------------------------
\section{Conclusion}
% -----------------------------------------------------------------------

Two Genetic Programming approaches were implemented and evaluated on the Breast Cancer Wisconsin dataset. The Arithmetic GP evolves symbolic arithmetic expressions, while the Decision Tree GP evolves binary decision trees using threshold comparisons. Both systems use ramped half-and-half initialisation, tournament selection, subtree crossover, and point mutation. Results from 30 independent seeded runs were compared using a Wilcoxon signed-rank test to assess statistical significance.

Two Genetic Programming approaches were implemented and evaluated on the Breast Cancer Wisconsin dataset. The Arithmetic GP evolves symbolic arithmetic expressions classified by the sign of the output, while the Decision Tree GP evolves binary decision trees using threshold comparisons on individual features. Both systems use ramped half-and-half initialisation, tournament selection, subtree crossover, and point mutation with a population of 200 over up to 100 generations.

Over 30 independent seeded runs, the Decision Tree GP achieved a higher and more stable mean test accuracy (50.50\% $\pm$ 3.65\%) compared to the Arithmetic GP (47.48\% $\pm$ 5.27\%). A Wilcoxon signed-rank test confirmed this difference is statistically significant ($p = 0.0088$). The best individual run overall belonged to the Arithmetic GP at 62.7907\% test accuracy (seed 25), suggesting that while the Arithmetic GP has higher potential, it is more sensitive to the random seed and less reliable across runs.

\end{document}
